\subsection{Construction corrélée et ordre des lecteurs}
\label{gen:orden-constructor-correlacionado}

L'antécédent commun des trois régions est l'émission arithmétique avec son registre de trajectoire. APP fournit les évaluations somme et produit ; TRIT apporte la signature orientée ; la transition \eqref{eq:actualizacion} compose sélection du mode, transport cardinal et mise à jour de la mémoire. Le recensement de cette émission produit le catalogue et ses fibres. Les prédicats de multiplicité, de support diagonal et d'orientation sélectionnent les représentants de \eqref{eq:palabras-regionales}. La corrélation régionale conserve cette provenance conjointe, ainsi que les incidences et les restrictions entre ses registres.

Les transports \(L_0,L_1\) agissent sur les représentants sélectionnés
dans la coordinatisation en lignes déclarée. Le calendrier et la concaténation distinguent
la longueur du préfixe, la phase visible et la mémoire accumulée. La
coordonnée profonde et ses quotients sont explicités ci-après par
\eqref{memprof:division} et \eqref{memprof:inversa}. Les restrictions des
historiques sont conservées au sens de \ref{hist:seccion}.
La sélection orientée précise de \(R_{36}\) possède sa démonstration
d'incidence en \ref{mg01:frontera}, après la construction de \(A_W\) ;
cet emplacement conserve les antécédents utilisés dans sa preuve.

Chaque lecteur régional ajoute une opération typée au registre qu'il reçoit.
Dans la clôture, le lecteur symétrique de la pentafibre fixe la pente
primitive et la loi de composition détermine le compensateur de
\eqref{eq:lector-pi}. Dans la propagation, la composition formelle, l'unité
et le générateur orienté déterminent la récurrence
\eqref{eq:propagacion}. Dans l'auto-échelle, le calendrier des modes produit
l'incidence \eqref{eq:fav}, dont la demi-droite positive fixe la coordonnée
normalisée. Les démonstrations suivantes conservent les normalisations
et les domaines de ces trois opérations.

Les bornes rationnelles et la séparation des cylindres évaluent ensuite les caractères obtenus. L'émission positionnelle effective et sa compatibilité sont démontrées dans \ref{act21:thm:df-emision-correcta} et \ref{act21:prop:df-prefijos-compatibles}. Le registre dodécaphasique conserve sa construction propre dans \ref{act21:chap:despliegue-selector-k} ; ses coordonnées et les trois lectures régionales interviennent ultérieurement dans la clôture entière de \ref{sec:alpha}. Ainsi, sélection des régions, transport, construction des caractères et évaluation des préfixes maintiennent leurs domaines et leur ordre, au sein de la même généalogie APP--TRIT--TPK.
